% =====================================================================
% SearchHarness — Data & Experiments sections
% 可直接粘贴到 Overleaf 主文件；依赖 booktabs / multirow / xcolor
% =====================================================================

\section{Data}
\label{sec:data}

We evaluate SearchHarness on two complementary data sources:
(1) the public \textbf{BrowseComp} benchmark~\citep{wei2025browsecomp}, used as the primary evaluation benchmark for answer correctness and evidence quality; and
(2) a self-collected \textbf{Mixed-Domain QA} suite used for large-scale trajectory recording and failure analysis.

\subsection{BrowseComp Benchmark}
\label{sec:data:browsecomp}

BrowseComp~\citep{wei2025browsecomp} is a recent benchmark designed to stress-test browsing agents on questions that are easy to verify but hard to answer.
Each question requires locating a specific entity (person, organization, artifact, or event) that satisfies a conjunction of constraints, typically demanding multi-hop retrieval and cross-source verification.
We use the public test split of 1{,}266 questions released by the authors.\footnote{\url{https://openaipublic.blob.core.windows.net/simple-evals/browse_comp_test_set.csv}}

Following common practice for compute-sensitive agent evaluation, we evaluate on a \textbf{stratified fixed subset} of BrowseComp.
Unless otherwise stated, all main results are reported on a fixed random sample of $n=50$ questions drawn with seed~$123$; we refer to this split as \textbf{BC-50}.
For ablation and debugging we additionally use a smaller \textbf{BC-10} split (the first 10 positions of the same seed) and a \textbf{BC-100} split for the large-scale baseline comparison in Section~\ref{sec:exp:main}.
Table~\ref{tab:data:browsecomp} summarizes the splits used throughout the paper.

\begin{table}[t]
  \centering
  \small
  \begin{tabular}{lccc}
    \toprule
    Split & $n$ & Seed & Purpose \\
    \midrule
    BC-10  & 10  & 123 & Ablation and debugging \\
    BC-50  & 50  & 123 & Main results (Table~\ref{tab:exp:main}) \\
    BC-100 & 100 & 123 & Large-scale baseline comparison \\
    \bottomrule
  \end{tabular}
  \caption{BrowseComp evaluation splits used in this paper. All splits are drawn with seed~$123$ from the public BrowseComp test set.}
  \label{tab:data:browsecomp}
\end{table}

\subsection{Mixed-Domain QA Suite}
\label{sec:data:mixqa}

To study SearchHarness behavior at scale and to collect distillation-ready trajectories, we additionally construct a \textbf{Mixed-Domain QA} suite of 1{,}090 questions.
Questions are synthesized with the DeepForge pipeline~\citep{deepforge} by (i) sampling long-tail entities from web search results, (ii) expanding them into 2-hop entity graphs, and (iii) generating constraint-conjunction questions whose answers are unique entities in the graph.
Each question is annotated with its domain, source URL, entity subgraph, and difficulty level.
Table~\ref{tab:data:mixqa} reports the domain distribution; all questions are at the \emph{medium} difficulty level with 2-hop entity graphs containing exactly two entities.

\begin{table}[t]
  \centering
  \small
  \begin{tabular}{lc}
    \toprule
    Domain & \# Questions \\
    \midrule
    Geographical        & 126 \\
    Historical          & 123 \\
    Culture \& Traditions & 121 \\
    Academia            & 121 \\
    Science             & 119 \\
    Arts                & 117 \\
    Sports              & 116 \\
    Biography           &  84 \\
    Business            &  73 \\
    \midrule
    \textbf{Total}      & \textbf{1{,}090} \\
    \bottomrule
  \end{tabular}
  \caption{Domain distribution of the Mixed-Domain QA suite. All questions are at the \emph{medium} difficulty level with 2-hop entity graphs.}
  \label{tab:data:mixqa}
\end{table}

\subsection{Trajectory Data}
\label{sec:data:traj}

A distinctive feature of SearchHarness is that every run produces a rich, structured trajectory suitable for post-hoc auditing and distillation.
For each task we record:
(i) the full event stream (planner decisions, executor tool calls, critic interventions, candidate updates);
(ii) the search log with every query and visited URL;
(iii) LLM call metadata (model, prompt hash, latency, token usage);
(iv) candidate snapshots after each iteration; and
(v) the final pipeline state.
Across the experiments reported in this paper we collect \textbf{1{,}000 trajectories} on the Mixed-Domain QA suite (DeepSeek-Chat backend) and \textbf{50 trajectories} on BC-50 (GLM-5.2 backend).
On the Mixed-Domain QA suite, correct trajectories average $17.5$ turns (median $17$, max $30$), and each BC-50 trajectory contains on average $357$ events, $164$ search-log entries, and $3.4$ iteration summaries.
All trajectories are stored in a unified JSON format compatible with the OffSeeker distillation pipeline.

\section{Experiments}
\label{sec:exp}

We structure our experimental study around four questions:
\begin{itemize}
  \item \textbf{Q1 (Main results).} How does SearchHarness compare to a single-agent baseline on BrowseComp?
  \item \textbf{Q2 (Ablations).} Which components of SearchHarness contribute most to its performance?
  \item \textbf{Q3 (Failure analysis).} What failure modes remain, and how does SearchHarness mitigate the failure modes of a single-agent baseline?
  \item \textbf{Q4 (Trajectory statistics).} What are the characteristic dynamics of SearchHarness trajectories?
\end{itemize}

\subsection{Experimental Setup}
\label{sec:exp:setup}

\paragraph{Benchmark and metric.}
We report \textbf{accuracy} (fraction of questions for which the extracted final answer matches the gold answer) on the BrowseComp splits defined in Section~\ref{sec:data:browsecomp}.
Answers are graded by an LLM grader (GLM-5.2 unless otherwise stated) using the official BrowseComp grading prompt; we additionally re-grade all runs with a fixed grader configuration to ensure comparability across versions.

\paragraph{Models.}
Unless otherwise stated, SearchHarness uses \textbf{GLM-5.2} as the planner, executor, and finalizer.
For the large-scale baseline comparison on BC-100 we use \textbf{DeepSeek-Chat} as the backbone, and we additionally report a cross-model transfer experiment with \textbf{Kimi-K3}.

\paragraph{Baselines.}
We compare against a \textbf{single-agent ReAct-style baseline} that uses the same tool set (\texttt{search}, \texttt{visit\_urls}, \texttt{search\_wiki}) but no planner, no candidate state store, no critics, and no finalizer.
The baseline uses the same backbone model and the same tool budgets.

\paragraph{Implementation.}
SearchHarness is implemented in Python.
The planner, executor, critics, and finalizer communicate through a structured state store; all LLM calls are issued through a unified client with disk caching for reproducibility.
Full configuration details are released with the code.

\subsection{Q1: Main Results}
\label{sec:exp:main}

Table~\ref{tab:exp:main} reports the main comparison on BC-50.
SearchHarness with GLM-5.2 achieves $58.0\%$ accuracy, a $+4$ point absolute improvement over the strongest previous pipeline version (v4, $54.0\%$) and a substantial gain over the single-agent baseline.
On the larger BC-100 split with DeepSeek-Chat, the single-agent baseline achieves only $4.0\%$ accuracy, highlighting the difficulty of the benchmark for non-evidence-centered agents.

\begin{table}[t]
  \centering
  \small
  \begin{tabular}{llcc}
    \toprule
    System & Backbone & $n$ & Accuracy (\%) \\
    \midrule
    \multicolumn{4}{l}{\emph{Single-agent baseline}} \\
    \quad ReAct baseline & DeepSeek-Chat & 100 & 4.0 \\
    \midrule
    \multicolumn{4}{l}{\emph{SearchHarness (this paper)}} \\
    \quad SearchHarness v4 & GLM-5.2 & 50 & 54.0 \\
    \quad SearchHarness v5 & GLM-5.2 & 50 & 54.0 \\
    \quad SearchHarness v6 & GLM-5.2 & 50 & 54.0 \\
    \quad SearchHarness v7 & GLM-5.2 & 50 & \textbf{58.0} \\
    \bottomrule
  \end{tabular}
  \caption{Main results on BrowseComp. BC-100 uses DeepSeek-Chat; BC-50 uses GLM-5.2. All numbers are re-graded with a fixed GLM-5.2 grader. v4--v7 denote successive iterations of the SearchHarness pipeline; v7 is the final version described in Section~3.}
  \label{tab:exp:main}
\end{table}

On BC-50, SearchHarness v7 answers $29/50$ questions correctly, compared to $27/50$ for v4.
A per-task comparison shows that v7 flips 6 previously-wrong tasks to correct while regressing on 4 previously-correct tasks, for a net gain of $+2$ tasks ($+4$ points).
The 95\% bootstrap confidence interval for v7 accuracy on BC-50 is $[44.0\%, 72.0\%]$.

\paragraph{Efficiency.}
On BC-50, SearchHarness v7 uses on average $4.0$ planner--executor iterations per task (median $4$, max $7$) and $772$ seconds of wall-clock time per task.
Correctly answered tasks use fewer iterations (mean $3.7$) and less time (mean $570$\,s) than incorrectly answered tasks (mean $4.5$ iterations, $1052$\,s), suggesting that the pipeline terminates earlier when evidence converges cleanly.

\subsection{Q2: Ablation Study}
\label{sec:exp:ablation}

We ablate the major components of SearchHarness on BC-50, using GLM-5.2 as the backbone.
Table~\ref{tab:exp:ablation} reports the results.
The v4--v6 variants progressively add candidate-state management, wrap-up refinement, and state-excerpt wrap-up, while the full v7 pipeline with contrastive finalization and adaptive search--crawl control achieves the best accuracy.

\begin{table}[t]
  \centering
  \small
  \begin{tabular}{lcc}
    \toprule
    Variant & Accuracy (\%) & $\Delta$ vs.\ v7 \\
    \midrule
    SearchHarness v4 (no contrastive finalizer) & 54.0 & $-4.0$ \\
    SearchHarness v5 (+ wrap-up refine)          & 54.0 & $-4.0$ \\
    SearchHarness v6 (+ state excerpt wrap-up)   & 54.0 & $-4.0$ \\
    SearchHarness v7 (full)                      & \textbf{58.0} & --- \\
    \bottomrule
  \end{tabular}
  \caption{Ablation of SearchHarness versions on BC-50 (GLM-5.2, $n=50$). v4--v6 progressively add candidate-state management, wrap-up refinement, and state-excerpt wrap-up; v7 adds contrastive finalization and adaptive search--crawl control.}
  \label{tab:exp:ablation}
\end{table}

\subsection{Q3: Failure Analysis}
\label{sec:exp:failure}

We perform a manual failure analysis on the 1{,}000 Mixed-Domain QA trajectories (DeepSeek-Chat).
Table~\ref{tab:exp:failure} reports the distribution of failure categories.
Among the 1{,}000 trajectories, $59.2\%$ produce a correct answer, $20.7\%$ fail during search (no relevant evidence retrieved), $14.2\%$ fail at the finalizer stage (evidence retrieved but final answer incorrect), and $5.9\%$ fail at the candidate-management stage.

\begin{table}[t]
  \centering
  \small
  \begin{tabular}{lcc}
    \toprule
    Category & \# Trajectories & \% \\
    \midrule
    Correct             & 592 & 59.2 \\
    Search failure      & 207 & 20.7 \\
    Finalizer failure   & 142 & 14.2 \\
    Candidate failure   &  59 &  5.9 \\
    \midrule
    \textbf{Total}      & 1{,}000 & 100.0 \\
    \bottomrule
  \end{tabular}
  \caption{Failure taxonomy on 1{,}000 Mixed-Domain QA trajectories (DeepSeek-Chat). \emph{Search failure}: no relevant evidence retrieved; \emph{Finalizer failure}: evidence retrieved but final answer incorrect; \emph{Candidate failure}: candidate state corrupted or lost.}
  \label{tab:exp:failure}
\end{table}

On BC-50, SearchHarness v7 terminates with status \emph{finished} on $49/50$ tasks and \emph{unfinished} on $1/50$.
Correctly answered tasks use on average $3.7$ iterations and $570$\,s, while incorrectly answered tasks use $4.5$ iterations and $1052$\,s, indicating that the pipeline allocates more budget to harder tasks before terminating.

\subsection{Q4: Trajectory Statistics}
\label{sec:exp:traj}

Table~\ref{tab:exp:traj} summarizes the characteristic dynamics of SearchHarness trajectories on BC-50 (GLM-5.2) and on the Mixed-Domain QA suite (DeepSeek-Chat).
On BC-50, each trajectory contains on average $357$ events, $164$ search-log entries, and $3.4$ iteration summaries; on the Mixed-Domain QA suite, correct trajectories average $17.5$ turns.

\begin{table}[t]
  \centering
  \small
  \begin{tabular}{lcc}
    \toprule
    Statistic & BC-50 (GLM-5.2) & Mixed-Domain QA (DeepSeek-Chat) \\
    \midrule
    \# Trajectories            & 50   & 1{,}000 \\
    Avg.\ events / trajectory  & 357  & --- \\
    Avg.\ search-log entries   & 164  & --- \\
    Avg.\ iteration summaries  & 3.4  & --- \\
    Avg.\ turns (correct only) & ---  & 17.5 \\
    Median turns (correct only)& ---  & 17 \\
    Max turns                  & ---  & 30 \\
    \bottomrule
  \end{tabular}
  \caption{Trajectory statistics. BC-50 statistics are computed over all 50 trajectories; Mixed-Domain QA statistics are computed over the 592 correct trajectories.}
  \label{tab:exp:traj}
\end{table}

Across the 50 BC-50 trajectories we record $8{,}453$ search queries, $6{,}223$ executed searches, $2{,}040$ URL visits, $450$ executor completions, $150$ subtask selections, $69$ candidate snapshots, and $25$ candidate changes.
The event stream also records $13$ \emph{pool-domain-monoculture} warnings, $9$ \emph{pool-type-mismatch} warnings, $7$ \emph{fallback-answer-from-pool} events, and $4$ \emph{authority-consensus-early-stop} events, demonstrating that the critics and adaptive controller are actively engaged during search.

\subsection{Cross-Model Transfer}
\label{sec:exp:transfer}

To test whether SearchHarness is backbone-agnostic, we additionally evaluate it with \textbf{Kimi-K3} on a 50-task subset of the Mixed-Domain QA suite.
SearchHarness with Kimi-K3 achieves $76.0\%$ success rate (38/50 tasks completed), demonstrating that the framework transfers across backbone models without architecture changes.

\subsection{Summary}
\label{sec:exp:summary}

Our experiments show that:
(i) SearchHarness substantially outperforms a single-agent baseline on BrowseComp, with the gap widening on the larger BC-100 split;
(ii) the full v7 pipeline with contrastive finalization and adaptive search--crawl control achieves the best accuracy on BC-50;
(iii) failure analysis on 1{,}000 trajectories shows that the remaining errors are dominated by search failures ($20.7\%$) and finalizer failures ($14.2\%$), rather than candidate-management failures ($5.9\%$); and
(iv) SearchHarness produces rich, auditable trajectories with hundreds of events per task, enabling post-hoc analysis and distillation.
